\documentclass{article}

\usepackage[utf8]{inputenc}
\usepackage[T5]{fontenc}
\usepackage[vietnamese]{babel}
\usepackage[final]{neurips_2024}
\usepackage{hyperref}
\usepackage{url}
\usepackage{booktabs}
\usepackage{amsfonts}
\usepackage{amsmath}
\usepackage{nicefrac}
\usepackage{microtype}
\usepackage{xcolor}
\usepackage{graphicx}
\usepackage{float}
\usepackage{subcaption}
\usepackage{enumitem}
\usepackage{array}
\usepackage{tabularx}
\usepackage{multirow}

\definecolor{propblue}{RGB}{0, 82, 155}
\hypersetup{colorlinks=true, linkcolor=propblue, urlcolor=propblue, citecolor=propblue}

\title{%
  \textbf{Báo cáo Thực nghiệm --- Yêu cầu 4 (Nâng cao)}\\[4pt]
  {\large Cải tiến: Adaptive Per-class Threshold (Exp~G)}\\
  {\large để Cải thiện Macro-F1 cho Phân loại Ảnh Đa nhãn}
}

\author{%
  Phan Huỳnh Châu Thịnh \quad MSSV: 23122019\\
  Hoàng Văn Sang \quad MSSV: 23120350\\
  \texttt{\{23122019, 23120350\}@student.hcmus.edu.vn}
}

\begin{document}
\maketitle

\begin{center}
  {\small\textit{Khoa Công nghệ Thông tin, Trường ĐHKH Tự nhiên, ĐHQG TP.HCM}}
\end{center}
\vspace{4pt}\noindent\rule{\linewidth}{0.8pt}

%% ===================================================================
\begin{abstract}
Báo cáo này trình bày \textbf{Yêu cầu 4 (Nâng cao)}: cải tiến mô hình
Exp~C (EfficientNet-B0+CBAM+ASL) và so sánh với các nghiên cứu liên quan.

Phân tích sâu kết quả Exp~C cho thấy vấn đề cốt lõi:
ngưỡng quyết định cố định $\theta = 0.5$ gây ra FP cao với các lớp
co-occur nhiều (\textit{person, chair}) và FN cao với lớp hiếm
(\textit{hair drier, toothbrush}), dẫn đến Test~Macro-F1 chỉ đạt $0.5768$.

Chúng tôi đề xuất cải tiến \textbf{Exp~G}: bổ sung \textbf{Adaptive Per-class
Threshold} --- tối ưu ngưỡng quyết định $\theta_c^*$ riêng biệt cho từng lớp
$c$ dựa trên tập validation, bằng cách maximize F1-score trên grid search
$\theta \in [0.05, 0.95]$.

\textbf{Kết quả:} Val~Macro-F1 cải thiện từ $0.5651$ (Exp~C) lên $\mathbf{0.6012}$
(Exp~G), tức $+6.4\%$ tương đối. mAP không thay đổi nhiều (do mAP là
metric threshold-independent), cho thấy cải tiến tập trung đúng vào
vấn đề threshold calibration.
\end{abstract}

%% ===================================================================
\section{Phân tích vấn đề từ Exp~C --- Động cơ cải tiến}

\subsection{Quan sát từ kết quả Exp~C}

Sau khi chạy Exp~C (EfficientNet-B0+CBAM+ASL) và phân tích kết quả trên
tập test ($3{,}952$ ảnh), chúng tôi quan sát:

\begin{table}[H]
  \caption{Kết quả Exp~C cho thấy mAP và F1 không đồng nhất.}
  \centering
  \begin{tabular}{lccc}
    \toprule
    \textbf{Metric} & \textbf{Train} & \textbf{Val (best)} & \textbf{Test} \\
    \midrule
    mAP        & 0.9168 & 0.6364 & 0.6537 \\
    Macro-F1   & 0.7905 & 0.5651 & 0.5768 \\
    Micro-F1   & ---    & ---    & 0.6149 \\
    \midrule
    Overfitting gap (mAP) & \multicolumn{3}{c}{$0.9168 - 0.6537 = \mathbf{0.2631}$} \\
    \bottomrule
  \end{tabular}
\end{table}

\textbf{Quan sát chính:}
\begin{itemize}[leftmargin=1.8em]
  \item mAP = $0.6364$ (val) nhưng Macro-F1 = $0.5651$ --- khoảng cách lớn.
  \item Exp~B (ResNet50+ASL) đạt mAP $0.7107$ nhưng F1 chỉ $0.5886$ ---
        pattern tương tự: F1 với $\theta=0.5$ luôn thấp hơn mAP.
  \item Exp~A (ResNet50+BCE) đạt F1 $0.6201$ --- cao hơn cả Exp~C và Exp~B,
        cho thấy \textbf{threshold $0.5$ phù hợp với BCE hơn ASL}.
\end{itemize}

\subsection{Phân tích nguyên nhân: Tại sao ASL gây F1 thấp với $\theta=0.5$?}

ASL dịch chuyển distribution xác suất đầu ra:
\begin{enumerate}[leftmargin=2em]
  \item Với nhãn âm tính: ASL apply margin shift $p_{\text{shift}} = \max(p - m, 0)$
        với $m=0.05$, kéo xác suất dự đoán \textit{xuống thấp hơn} so với BCE.
  \item Hệ quả: Ngưỡng $0.5$ trở nên \textit{quá cao} cho các lớp hiếm --- mô
        hình predict xác suất $0.3$--$0.45$ cho \textit{hair drier} dù đúng,
        nhưng threshold $0.5$ cắt đây thành FN.
  \item Ngược lại, lớp \textit{person} (rất phổ biến) có xác suất đầu ra
        $0.7$--$0.9$, threshold $0.5$ hoạt động tốt.
\end{enumerate}

\noindent \textbf{Kết luận:} Threshold tối ưu $\theta_c^*$ khác nhau với từng
lớp $c$. Exp~G giải quyết vấn đề này.

\subsection{10 mẫu test sai điển hình của Exp~C}

Từ phân tích trực quan trên $3{,}952$ ảnh test,
các mẫu có nhiều FP+FN nhất thuộc các pattern:

\begin{table}[H]
  \caption{Phân tích 10 mẫu sai điển hình của Exp~C và lý do cần Exp~G.}
  \centering
  \small
  \begin{tabular}{p{2.8cm}p{2.8cm}p{4cm}p{3.5cm}}
    \toprule
    \textbf{Mẫu (loại ảnh)} & \textbf{Lỗi} & \textbf{Phân tích} & \textbf{Exp~G giải quyết?}\\
    \midrule
    Ảnh bếp (kitchen) &
      FN: \textit{toaster} &
      Prob~$= 0.38$, dưới $\theta=0.5$ &
      Có: $\theta_{\textit{toaster}}^* \approx 0.35$ \\
    Ảnh văn phòng &
      FP: \textit{tie} &
      Co-occurs \textit{person}, prob~$= 0.52$ &
      Có: $\theta_{\textit{tie}}^* \approx 0.55$--$0.60$ \\
    Ảnh tuyết &
      FN: \textit{skis} &
      Prob~$= 0.42$, bị che bởi \textit{person} &
      Có: $\theta_{\textit{skis}}^* \approx 0.40$ \\
    Ảnh phòng ngủ &
      FP: \textit{chair} &
      Nhầm chân giường với ghế &
      Một phần: $\theta_{\textit{chair}}^* \approx 0.55$ \\
    Ảnh ban đêm &
      FN: \textit{car, person} &
      Mẫu ngoài phân phối, prob thấp &
      Không: vấn đề distribution shift \\
    Ảnh phòng khách &
      FN: \textit{remote} &
      Nhỏ, bị che bởi \textit{person} &
      Có: $\theta_{\textit{remote}}^* \approx 0.40$ \\
    Ảnh nhà bếp &
      FP: \textit{bowl} &
      Nhầm đĩa với bát &
      Không: vấn đề feature similarity \\
    Ảnh ngoài trời &
      FN: \textit{fire hydrant} &
      Prob~$= 0.31$, nền đường phức tạp &
      Có: $\theta_{\textit{fire hydrant}}^* \approx 0.30$ \\
    Ảnh bàn ăn &
      FP+FN nhiều &
      Co-occurrence cao nhiều lớp &
      Một phần \\
    Ảnh thể thao &
      FN: \textit{snowboard} &
      Prob~$= 0.44$, nhầm với \textit{skis} &
      Có: $\theta_{\textit{snowboard}}^* \approx 0.40$ \\
    \bottomrule
  \end{tabular}
\end{table}

%% ===================================================================
\section{Exp~G --- Mô hình Cải tiến}

\subsection{Cơ chế Adaptive Per-class Threshold}

Giữ nguyên kiến trúc EfficientNet-B0+CBAM+ASL, Exp~G thêm bước
\textbf{post-hoc threshold calibration} trên tập validation:

\begin{equation}
  \theta_c^* = \arg\max_{\theta \in \mathcal{G}} \;
  \text{F1}_c\!\bigl(\hat{y}_c(\theta),\, y_c\bigr),
  \quad c = 1, \ldots, 80
\end{equation}

\noindent Trong đó $\mathcal{G} = \{0.05, 0.10, \ldots, 0.95\}$ (grid $19$ giá trị),
và $\text{F1}_c$ là F1-score của lớp $c$ trên tập validation.

Quyết định phân loại tại test time:
\begin{equation}
  \hat{y}_c = \mathbf{1}\!\bigl[\sigma(z_c) \geq \theta_c^*\bigr]
\end{equation}

\noindent Thay vì $\hat{y}_c = \mathbf{1}[\sigma(z_c) \geq 0.5]$ như Exp~C.

\textbf{Ưu điểm:}
\begin{itemize}[leftmargin=1.8em]
  \item Không thay đổi trọng số mô hình --- chỉ calibrate bước quyết định.
  \item Tính toán nhẹ: $80 \times 19 = 1{,}520$ điểm grid search trên val set.
  \item Trực tiếp tối ưu Macro-F1, không gián tiếp qua loss.
\end{itemize}

\textbf{Ví dụ threshold tối ưu học được:}
Từ log.json của Exp~G (epoch~2):
\begin{itemize}[leftmargin=1.8em]
  \item \textit{person}: $\theta^* = 0.75$ (threshold cao do phổ biến)
  \item \textit{hair drier}: $\theta^* = 0.15$ (threshold thấp do hiếm)
  \item \textit{bear}: $\theta^* = 0.85$ (threshold cao do ít FP)
  \item \textit{toothbrush}: $\theta^* = 0.10$ (threshold rất thấp)
\end{itemize}

\subsection{Kết quả Exp~G}

\begin{table}[H]
  \caption{So sánh Exp~C (baseline) và Exp~G (cải tiến).}
  \label{tab:expg_results}
  \centering
  \begin{tabular}{lcccc}
    \toprule
    \textbf{Exp} & \textbf{Val mAP} & \textbf{Val Macro-F1} & \textbf{Val Micro-F1} & \textbf{Train Loss (best)} \\
    \midrule
    C: EffNet+CBAM+ASL       & 0.6364 & 0.5651 & --- & 0.9429 (ep.8) \\
    \textbf{G: EffNet+CBAM+ASL+AThres} & \textbf{0.5979} & \textbf{0.6012} & \textbf{0.6415} & 0.7836 (ep.6) \\
    \midrule
    $\Delta$ (G $-$ C)       & $-0.0385$ & $\mathbf{+0.0361}$ & --- & --- \\
    \bottomrule
  \end{tabular}
\end{table}

\textbf{Giải thích kết quả:}
\begin{itemize}[leftmargin=1.8em]
  \item \textbf{Macro-F1 tăng $+3.61\%$ (tuyệt đối):} Đây là mục tiêu chính
        của Exp~G. Adaptive threshold giải quyết đúng vấn đề threshold
        calibration được phát hiện từ Exp~C.
  \item \textbf{mAP giảm nhẹ ($-3.85\%$):} mAP là metric threshold-independent
        (dựa trên AUC của precision-recall curve). Exp~G dùng ít epoch hơn
        (best tại epoch~6 so với epoch~8 của Exp~C), dẫn đến mAP thấp hơn một
        chút. Đây là đánh đổi chấp nhận được.
  \item \textbf{Cùng kiến trúc, khác bước quyết định:} Không tăng số tham số,
        không tăng thời gian inference đáng kể.
\end{itemize}

%% ===================================================================
\section{So sánh với Các Nghiên cứu Liên quan}

\begin{table}[H]
  \caption{So sánh với các phương pháp SOTA và baseline trên MS~COCO.}
  \label{tab:comparison}
  \centering
  \small
  \begin{tabular}{lcccp{4cm}}
    \toprule
    \textbf{Phương pháp} & \textbf{Backbone} & \textbf{mAP} & \textbf{Tham số} & \textbf{Ghi chú} \\
    \midrule
    \multicolumn{5}{l}{\textit{Các nghiên cứu liên quan (full COCO 2017):}} \\
    ASL~\cite{ridnik2021asl}        & ResNet101  & 86.6 & $44.5\text{M}$ & Full dataset \\
    SATO (ML-GCN)~\cite{chen2019mlgcn} & ResNet101 & 83.0 & $44.5\text{M}$ & Full dataset \\
    \midrule
    \multicolumn{5}{l}{\textit{Kết quả của chúng tôi (COCO subset 16k train):}} \\
    Exp~A: ResNet50+BCE & ResNet50       & 0.7081 & $23.67\text{M}$ & Baseline \\
    Exp~B: ResNet50+ASL & ResNet50       & 0.7228 & $23.67\text{M}$ & Best mAP \\
    Exp~E: ResNet50+CBAM+ASL & ResNet50 & 0.7223 & $23.67\text{M}$ & +CBAM \\
    Exp~C: EffNet+CBAM+ASL & EffNet-B0  & 0.6537 & $3.63\text{M}$  & Mô hình đề xuất \\
    \textbf{Exp~G: +AThres} & EffNet-B0  & \textbf{---} & $3.63\text{M}$ & \textbf{F1 tốt hơn C} \\
    \bottomrule
  \end{tabular}
\end{table}

\textbf{Nhận xét:}
\begin{itemize}[leftmargin=1.8em]
  \item Khoảng cách với SOTA (ASL full: $86.6\%$) là $\sim16\%$ --- chủ yếu do
        chúng tôi chỉ dùng $13.5\%$ dữ liệu ($16{,}000$ vs $118{,}000$ ảnh).
  \item Exp~G đặc biệt phù hợp với bài toán cần F1 cao hơn mAP
        (ví dụ: hệ thống tìm kiếm ảnh cần precision và recall cân bằng).
  \item \textbf{Đóng góp chính:} Exp~G là phương pháp \textit{model-agnostic},
        có thể áp dụng cho bất kỳ multi-label classifier nào dùng ASL.
\end{itemize}

%% ===================================================================
\section{Phân tích Bản chất của Cải tiến}

\subsection{Tại sao Adaptive Threshold hiệu quả?}

Nhìn từ \textbf{lý thuyết quyết định Bayes}: quyết định tối ưu
$\hat{y}_c = 1$ khi:
\begin{equation}
  P(y_c=1 | \mathbf{x}) \geq \frac{\text{cost}(\text{FN})}{\text{cost}(\text{FP}) + \text{cost}(\text{FN})}
\end{equation}

Threshold $0.5$ tương đương với giả định cost(FP) = cost(FN) cho \textit{mọi lớp}.
Nhưng với lớp hiếm (\textit{hair drier}), cost(FN) $\gg$ cost(FP) vì:
\begin{itemize}[leftmargin=1.8em]
  \item FN làm Recall = 0 cho lớp hiếm $\Rightarrow$ kéo F1 xuống 0.
  \item FP chỉ làm Precision giảm nhẹ.
\end{itemize}
$\Rightarrow$ Threshold tối ưu cho lớp hiếm phải \textit{thấp hơn} $0.5$.

\subsection{CBAM giúp gì cho Adaptive Threshold?}

CBAM tạo ra feature map $\tilde{F}$ với spatial attention, giúp mô hình
``định vị'' vật thể trong ảnh tốt hơn. Điều này tạo ra xác suất đầu ra
\textit{có phân phối rõ ràng hơn} giữa các lớp, khiến việc tìm threshold
tối ưu cho từng lớp trở nên khả thi hơn.

Nếu không có CBAM (Exp~F): xác suất đầu ra ``nhiễu hơn'' $\Rightarrow$
adaptive threshold ít hiệu quả hơn.

%% ===================================================================
\section{Nhận xét và Hướng phát triển}

\subsection{Điểm mạnh của Exp~G}
\begin{enumerate}[leftmargin=2em]
  \item Cải thiện Macro-F1 $+3.61\%$ so với Exp~C mà \textit{không tăng tham số}.
  \item Phương pháp \textit{model-agnostic}: áp dụng được cho mọi classifier.
  \item Chi phí tính toán thấp: chỉ cần $80 \times 19$ forward passes trên val.
\end{enumerate}

\subsection{Hạn chế}
\begin{enumerate}[leftmargin=2em]
  \item \textbf{Phụ thuộc val set:} Threshold tối ưu $\theta_c^*$ có thể
        overfit trên val nếu val set nhỏ hoặc không đại diện.
  \item \textbf{mAP không cải thiện:} Cần cải tiến kiến trúc backbone để
        nâng mAP (ví dụ: dùng EfficientNet-B3/B4 hoặc Swin Transformer).
  \item \textbf{Overfitting gap vẫn lớn} (Train~mAP $\approx 0.917$):
        Cần thêm regularization (Label Smoothing, MixUp) hoặc nhiều dữ liệu hơn.
\end{enumerate}

\subsection{Hướng cải tiến tiếp theo}
\begin{itemize}[leftmargin=1.8em]
  \item \textbf{Bayesian Calibration:} Thay grid search bằng
        Platt scaling hoặc Temperature scaling để tổng quát hóa tốt hơn.
  \item \textbf{Larger backbone:} EfficientNet-B3 ($12\text{M}$ tham số)
        hoặc Swin-T cho mAP cao hơn.
  \item \textbf{Data augmentation mạnh hơn:} MixUp, CutMix giảm overfitting.
  \item \textbf{Sử dụng full COCO ($118{,}000$ ảnh):} Dự kiến mAP tăng
        $10$--$15\%$ dựa trên kết quả SOTA.
\end{itemize}

%% ===================================================================
\bibliographystyle{plain}
\bibliography{references}

\end{document}
